\section*{अध्याय \ref{ch:g8:ice-on-earth} --- पृथ्वी पर बर्फ़}

\begin{solution}{exo:g8:ice-on-earth:1}
(क) संचयन क्षेत्र। (ख) अपक्षरण क्षेत्र। (ग) मुख। (घ) एक (अंतस्थ)
\omterm{def:g8:ice-on-earth:glacial-landforms}{हिमोढ़}।
\end{solution}

\begin{solution}{exo:g8:ice-on-earth:2}
(क) ग़लत: \omterm{def:g8:ice-on-earth:glacier}{हिमनद} ज़मीन पर गिरे और बर्फ़ में बदले \omterm{def:g4:clouds-rain-snow:precipitation}{हिम} का बना होता है; जमा हुआ समुद्री
पानी \omterm{def:g8:ice-on-earth:sea-ice}{समुद्री बर्फ़} है। (ख) सही। (ग) ग़लत: उसकी बर्फ़ अब भी ढलान से नीचे बहती है; मुख
इसलिए पीछे हटता है क्योंकि बर्फ़ जितना बहती है उससे तेज़ पिघलती है। (घ) सही।
\end{solution}

\begin{solution}{exo:g8:ice-on-earth:3}
(क) \omterm{def:g8:ice-on-earth:glacial-landforms}{U-आकार की घाटी}। (ख) अनियत शिला। (ग) \omterm{def:g8:ice-on-earth:ice-sheet}{हिम शेल्फ़}। (घ) \omterm{def:g8:ice-on-earth:permafrost}{पर्माफ़्रॉस्ट}।
\end{solution}

\begin{solution}{exo:g8:ice-on-earth:4}
(क) अंटार्कटिका और ग्रीनलैंड। (ख) अंटार्कटिका, लगभग \qty{58}{m} (ग्रीनलैंड: लगभग
\qty{7.4}{m})। (ग) लगभग दसवाँ हिस्सा (10~\%)।
\end{solution}

\begin{solution}{exo:g8:ice-on-earth:5}
(क) 1979 में लगभग \num{7.05}~मिलियन~\unit{km^2} और 2012 में
\num{3.57}~मिलियन~\unit{km^2}। (ख) $7.05 - 3.57 = 3.48$: लगभग
\num{3.5}~मिलियन~\unit{km^2}। (ग) $3.48 \div 7.05 \approx 0.49$: लगभग आधा। (घ) कोई
अकेला वर्ष उस गर्मी के मौसम पर निर्भर करता है: 2012 रेखा से बहुत नीचे है। प्रवृत्ति रेखा
सभी वर्षों का उपयोग करती है, इसलिए एक असामान्य गर्मी उसे बहुत कम बदलती है।
\end{solution}

\begin{solution}{exo:g8:ice-on-earth:6}
(क) $2004 - 1941 = 63$ वर्ष। (ख) $12000 \div 63 \approx 190.5$: लगभग
\qty{190}{m} प्रति वर्ष। (ग) $2021 - 1984 = 37$ वर्ष; $3200 \div 37 \approx
86$:
लगभग \qty{86}{m} प्रति वर्ष। (घ) म्यूर \omterm{def:g8:ice-on-earth:glacier}{हिमनद}, दोगुने से भी ज़्यादा तेज़।
\end{solution}

\begin{solution}{exo:g8:ice-on-earth:7}
\emph{संचयन क्षेत्र} में, पहाड़ पर ऊँचे, हर वर्ष जितना \omterm{def:g4:clouds-rain-snow:precipitation}{हिम} पिघलता है उससे ज़्यादा गिरता है।
नया \omterm{def:g4:clouds-rain-snow:precipitation}{हिम} पुराने \omterm{def:g4:clouds-rain-snow:precipitation}{हिम} को दबाता और ढकता है, जो लगभग एक वर्ष बाद \emph{फ़िर्न} बन जाता है,
फिर, और गहरे और बहुत बाद में, ठोस बर्फ़। अपने ही भार से बर्फ़ \omterm{def:g1:landscapes-around-us:landform}{घाटी} में धीरे-धीरे नीचे
\emph{बहती है}। और नीचे, \emph{अपक्षरण क्षेत्र} में, वह \omterm{def:g4:clouds-rain-snow:precipitation}{हिम} गिरने से ज़्यादा तेज़ी से
पिघलती है; मुख तक पहुँचने वाली बर्फ़ वह बर्फ़ है जो ऊपर, अक्सर बहुत पहले,
बनी थी।
\end{solution}

\begin{solution}{exo:g8:ice-on-earth:8}
(क) गिलास B में: पिघली बर्फ़ पानी में बहकर उसे ऊपर उठा देती है; गिलास A में स्तर जहाँ
था वहीं रहता है। (ख) गिलास A \omterm{def:g8:ice-on-earth:ice-sheet}{हिमशैल} है (\omterm{def:g2:water-on-earth:ocean}{समुद्र} में पहले से तैरती बर्फ़); गिलास B
\omterm{def:g8:ice-on-earth:ice-sheet}{हिमचादर} है (ज़मीन पर की बर्फ़)। (ग) \omterm{def:g8:ice-on-earth:ice-sheet}{हिम शेल्फ़} अपने पीछे के \omterm{def:g8:ice-on-earth:glacier}{हिमनदों} को रोके रखती है।
जब वह टूटती है, तो ये \omterm{def:g8:ice-on-earth:glacier}{हिमनद} तेज़ हो जाते हैं और ज़मीन की और बर्फ़ \omterm{def:g2:water-on-earth:ocean}{समुद्र} में उड़ेलते हैं,
और इससे \omterm{def:g2:water-on-earth:ocean}{समुद्र} ऊपर उठता है।
\end{solution}

\begin{solution}{exo:g8:ice-on-earth:9}
(क) \omterm{def:g8:ice-on-earth:glacial-landforms}{U-आकार की घाटी}, जिसे एक \omterm{def:g8:ice-on-earth:glacier}{हिमनद} ने तराशा। (ख) एक \omterm{def:g8:ice-on-earth:glacier}{हिमनद} उसे \omterm{def:g1:landscapes-around-us:landform}{घाटी} में बहुत ऊपर से,
जहाँ ग्रेनाइट है, लाया और बर्फ़ पिघलने पर छोड़ गया: यह एक अनियत शिला है। (ग) एक पार्श्व
\omterm{def:g8:ice-on-earth:glacial-landforms}{हिमोढ़}: उसकी ऊँचाई दिखाती है कि \omterm{def:g8:ice-on-earth:glacier}{हिमनद} कभी कितना मोटा था; बर्फ़ \omterm{def:g1:landscapes-around-us:landform}{घाटी} की ढलान पर कम से
कम उतनी ऊँचाई तक पहुँचती थी।
\end{solution}

\begin{solution}{exo:g8:ice-on-earth:10}
(क) $7.4 \div 10 = \qty{0.74}{m}$। (ख) $58 \div 100 = \qty{0.58}{m}$।
(ग) $58 \div 7.4 \approx 7.8$: लगभग 8 गुना ज़्यादा।
\end{solution}

\begin{solution}{exo:g8:ice-on-earth:11}
(क) $7.21 \times 10^6 = \qty{7210000}{km^2}$ और $4.65 \times 10^6 =
\qty{4650000}{km^2}$।
(ख) $7.21 - 4.65 = \num{2.56}\text{ मिलियन }\unit{km^2}$, यानी \qty{2560000}{km^2}।
(ग) $2.56 \div 3 \approx 0.85$: लगभग \num{0.85}~मिलियन~\unit{km^2} प्रति दशक,
प्रवृत्ति रेखा के \num{0.76}~मिलियन~\unit{km^2} से थोड़ा ज़्यादा; हर दशक का औसत अब भी
अपनी दस गर्मियों के उतार-चढ़ाव लिए हुए है, पर ग्राफ़ पढ़ने के दोनों तरीक़े एक जैसी बड़ी
कमी दिखाते हैं। (घ) विस्तार शून्य से नीचे नहीं गिर सकता, और प्रवृत्ति रेखा केवल मापे गए
वर्षों का वर्णन करती है: आगे क्या होगा, यह इस पर निर्भर है कि आर्कटिक की \omterm{def:g5:weather-and-climate:climate}{जलवायु} कैसे
बदलती है, इसलिए रेखा कोई पूर्वानुमान नहीं है।
\end{solution}

\begin{solution}{exo:g8:ice-on-earth:12}
(क) $120 \div 1 = 120$: \omterm{def:g1:landscapes-around-us:water}{तट} \omterm{def:g2:water-on-earth:ocean}{समुद्र} की ओर लगभग \qty{120}{km} और आगे था। (ख) वह
\omterm{def:g2:water-on-earth:ocean}{महासागर} से भाप बनकर उड़ा था, \omterm{def:g3:continents-and-oceans:continent}{महाद्वीपों} पर \omterm{def:g4:clouds-rain-snow:precipitation}{हिम} बनकर गिरा था और वहाँ बड़ी \omterm{def:g8:ice-on-earth:ice-sheet}{हिमचादरों} की
बर्फ़ के रूप में रुका रहा। (ग) अंतस्थ \omterm{def:g8:ice-on-earth:glacial-landforms}{हिमोढ़} \omterm{def:g8:ice-on-earth:glacier}{हिमनद} के छोर पर बनता है। किसी भी पहाड़ से दूर
किसी \omterm{def:g1:landscapes-around-us:landform}{मैदान} पर मिला \omterm{def:g8:ice-on-earth:glacial-landforms}{हिमोढ़} दिखाता है कि कभी कोई \omterm{def:g8:ice-on-earth:ice-sheet}{हिमचादर} उस \omterm{def:g1:landscapes-around-us:landform}{मैदान} तक पहुँची थी, उस समय
जब बर्फ़ आज से कहीं ज़्यादा ज़मीन ढकती थी।
\end{solution}

\begin{solution}{pb:g8:ice-on-earth:1}
\textbf{1.} $1080 - 0 = \qty{1080}{m}$।

\textbf{2.} 1960--1970: \qty{180}{m}; 1970--1980: \qty{150}{m};
1980--1990: \qty{190}{m}; 1990--2000: \qty{180}{m}; 2000--2010:
\qty{180}{m}; 2010--2020: \qty{200}{m}। सबसे कम 1970--1980 में, सबसे ज़्यादा
2010--2020 में।

\textbf{3.} नहीं: पीछे हटता \omterm{def:g8:ice-on-earth:glacier}{हिमनद} भी ढलान से नीचे बहता है। उसका मुख इसलिए पीछे
हटता है क्योंकि बर्फ़ जितनी तेज़ी से आती है, उससे तेज़ पिघलती है।

\textbf{4.} एक अंतस्थ \omterm{def:g8:ice-on-earth:glacial-landforms}{हिमोढ़}: जब मुख वहाँ खड़ा था, तब बर्फ़ चट्टानी मलबे को \omterm{def:g1:landscapes-around-us:landform}{घाटी} में नीचे
लाई और उसे अपने मुख पर, जहाँ बर्फ़ पिघलती थी, गिरा दिया।

\textbf{5.} $1080 \div 6 = \qty{180}{m}$ प्रति दशक।

\textbf{6.} $180 \div 10 = \qty{18}{m}$ प्रति वर्ष।

\textbf{7.} $190 \div 18 \approx 10.6$: लगभग 11 गुना तेज़।

\textbf{8.} कम: उसने अपने अपक्षरण क्षेत्र में जितनी बर्फ़ खोई, उससे कम अपने संचयन क्षेत्र
में पाई, इसलिए उसका मुख पीछे हटा।

\textbf{9.} U-आकार का: खड़ी ढलानों वाला चौड़ा तल, क्योंकि \omterm{def:g8:ice-on-earth:glacier}{हिमनद} ने \omterm{def:g1:landscapes-around-us:landform}{घाटी} को भरे रहते हुए
उसे चौड़ा और गहरा किया।

\textbf{10.} $1980 - 1080 = \qty{900}{m}$।

\textbf{11.} $900 \div 18 = 50$ वर्ष: \textbf{2070} में।

\textbf{12.} चाल दशक-दर-दशक बदली (150 से \qty{200}{m} तक), और भविष्य का पीछे हटना
आने वाले वर्षों के हिमपात और गर्मियों के तापमानों पर निर्भर है: बीती दर कोई पूर्वानुमान
नहीं है।
\end{solution}
